(1) 横向第 m 阶空间频率 \(k_{x,m}=2\pi m/d\)，总波数 \(k=2\pi/\lambda\)。近轴展开纵向波数
\[
k_{z,m}=\sqrt{k^2-k_{x,m}^2}\simeq k-\frac{k_{x,m}^2}{2k}
=k-\frac{\pi\lambda m^2}{d^2}.
\]
去掉全部谐波共有的 \(e^{ikz}\) 后，场包络为
\[
u(x,z)e^{-ikz}=\sum_{m\in\mathbb Z}c_m e^{i2\pi mx/d}
e^{-i\pi\lambda z m^2/d^2}.
\]
二次 m 相位是自成像的关键，不能只考虑一级衍射。

(2) 令 \(z_T=2d^2/\lambda\)，相对相位为 \(e^{-i2\pi m^2}=1\)，故
\[
u(x,z_T)=e^{ikz_T}u(x,0).
\]
所有空间傅里叶阶严格重相合，强度图样复现。此结论要求保留的衍射阶满足 \(k_{x,m}\ll k\)，否则精确波数的高阶校正会破坏理想 Talbot 复相。

(3) 在 z=zT/2=d²/λ，传播相位为 \(e^{-i\pi m^2}=(-1)^m=e^{i\pi m}\)（因为整数 \(m^2\) 与 m 奇偶性相同）。所以
\[
u(x,z_T/2)=e^{ikz_T/2}\sum_m c_m e^{i2\pi m(x+d/2)/d}
=e^{ikz_T/2}u(x+d/2,0).
\]
这表示整幅复场横向平移 d/2；强度若不是半周期对称，其图样相对于原光栅也发生同样平移。

(4) 在 z=zT/4 时相位为 \(e^{-i\pi m^2/2}\)：偶 m 取1、奇 m 取 \(-i\)。将初场拆为偶阶 \(u_e\) 和奇阶 \(u_o\)，有 \(u(z_T/4)e^{-ikz_T/4}=u_e-i u_o\)。又 \(u(x+d/2,0)=u_e-u_o\)，故
\[
u_e-i u_o=\frac{1-i}{2}u(x,0)+\frac{1+i}{2}u(x+d/2,0).
\]
此时是原场与半周期移位场的相干叠加，通常并不等于它们的强度直接相加，是分数 Talbot 自成像的基本机制。
